\input{LectureSlides/includes/metropolis_preamble}


% Title page info
\title[Ordinal Games]{Ordinal Games in Strategic Form}
\author[ECON 5001]{ECON 5001\\Prof. Healy} \color{metrop}
\institute[OSU]{}
\date[]{\vfill {\tiny Updated \today\ at\ \DTMcurrenttime}}

\begin{document}

\frame{\maketitle}

\begin{frame}{Where We're Headed}
\begin{itemize}
    \item A \textbf{game-frame}: who plays, what they can do, what happens
    \item \textbf{Preferences}: what players \textit{want}
    \item Frame $+$ preferences $=$ \textbf{game}
    \item \textbf{Utility} as a way to write preferences down
    \item \textbf{Dominance}: our first solution concept
\end{itemize}
\vfill
Big idea for today:
\bsk
\vspace{0.6in}
\end{frame}

% ==============================================================
% 1. GOLDEN BALLS
% ==============================================================

\begin{frame}{Golden Balls}
A British game show. Two contestants, Sarah and Steven, have built up a jackpot of \$100,000.
\bsk
Each secretly picks one of two balls:
\begin{itemize}
    \item one says \textbf{Split}
    \item one says \textbf{Steal}
\end{itemize}
They choose \textbf{simultaneously}. Then both are revealed.
\bsk
\begin{itemize}
    \item Both Split: \hspace{0.2in} each gets \$50,000
    \item One Steals: \hspace{0.14in} the stealer gets all \$100,000, the other gets nothing
    \item Both Steal: \hspace{0.19in} nobody gets anything
\end{itemize}
\vspace{-0.5in}
\end{frame}

\begin{frame}{Golden Balls: The Table}
\centering
\begin{tabular}{cc|c|c|}
\cline{3-4}
& & \multicolumn{2}{c|}{\textbf{Steven}} \\
\cline{3-4}
& & Split & Steal \\
\cline{2-4}
\multicolumn{1}{c|}{} & \multicolumn{1}{|c|}{Split}
  & \makecell{Sarah: \$50k\\Steven: \$50k}
  & \makecell{Sarah: \$0\\Steven: \$100k} \\
\cline{2-4}
\multicolumn{1}{c|}{\textbf{Sarah}} & \multicolumn{1}{|c|}{Steal}
  & \makecell{Sarah: \$100k\\Steven: \$0}
  & \makecell{Sarah: \$0\\Steven: \$0} \\
\cline{2-4}
\end{tabular}
\bsk

\raggedright
Each row is a choice for Sarah, each column a choice for Steven.\\
Each cell is what \textbf{happens}.
\vfill
So: what should Sarah do?
\vspace{-0.75in}
\end{frame}

\begin{frame}{A Tempting Argument}
\centering
{\footnotesize
\begin{tabular}{cc|c|c|}
\cline{3-4}
& & \multicolumn{2}{c|}{Steven} \\
\cline{3-4}
& & Split & Steal \\
\cline{2-4}
\multicolumn{1}{c|}{} & \multicolumn{1}{|c|}{Split} & \$50k,\ \$50k & \$0,\ \$100k \\
\cline{2-4}
\multicolumn{1}{c|}{Sarah} & \multicolumn{1}{|c|}{Steal} & \$100k,\ \$0 & \$0,\ \$0 \\
\cline{2-4}
\end{tabular}\\
\vspace{2pt}
{\scriptsize Sarah's amount listed first}}
\bsk

\raggedright
Sarah can't control Steven. But she can consider both cases.
\begin{itemize}
    \item Suppose Steven picks \textbf{Steal}: \vspace{0.45in}
    \item Suppose Steven picks \textbf{Split}: \vspace{0.45in}
    \item Therefore Sarah should \underline{\hspace{1.6in}}
\end{itemize}
\end{frame}

\begin{frame}{What's Wrong With That Argument?}
\vfill
\centering
\large
The argument snuck in an assumption.
\vfill
\normalsize
\raggedright
\vspace{1.2in}
\end{frame}

\begin{frame}{What's Wrong With That Argument?}
We assumed Sarah is \textbf{selfish and greedy}: she cares only about her own money, and more is better.
\bsk
That might be false! Sarah might:
\begin{itemize}
    \item value \textbf{fairness} --- an even split beats grabbing everything
    \item feel \textbf{benevolence} --- if she gets nothing anyway, she'd rather Steven get something
    \item feel \textbf{spite} --- she'd rather Steven get nothing
\end{itemize}
\bsk
Under some of these, the ``obvious'' answer flips completely.
\vfill
\textbf{Moral:} the table above is not yet enough to say what's rational.
\vspace{-0.4in}
\end{frame}

% ==============================================================
% 2. GAME FRAMES
% ==============================================================

\begin{frame}{Frames vs.\ Games}
What the table gave us:
\begin{itemize}
    \item the players
    \item what each can do
    \item what \textbf{outcome} results from each combination
\end{itemize}
\bsk
What the table did \textit{not} give us:
\begin{itemize}
    \item how each player \textbf{ranks} those outcomes
\end{itemize}
\vfill
\centering
\large
Table alone $=$ \textbf{game-frame}\\
\bsk
Frame $+$ rankings $=$ \textbf{game}
\vfill
\end{frame}

\begin{frame}{Our Own Example: Friday Night}
Alex and Blair want to have dinner. Both phones are dead, so no messaging.\\
Each just walks to one of the two places they always go.
\bsk
\begin{itemize}
    \item Alex chooses: Taco Truck ($T$) or Ramen Shop ($R$)
    \item Blair chooses: Taco Truck ($T$) or Ramen Shop ($R$)
    \item They choose simultaneously (neither sees the other's choice)
\end{itemize}
\bsk
Four things can happen:
\begin{itemize}
    \item $o_1$: both at the Taco Truck
    \item $o_2$: Alex at tacos, Blair at ramen (they eat alone)
    \item $o_3$: Alex at ramen, Blair at tacos (they eat alone)
    \item $o_4$: both at the Ramen Shop
\end{itemize}
\end{frame}

\begin{frame}{Friday Night as a Frame}
\begin{itemize}
    \item Players: \hspace{0.5in} $I=\{1,2\}$ \hspace{0.2in} (1 $=$ Alex, 2 $=$ Blair)
    \item Strategies: \hspace{0.28in} $S_1=\{T,R\}$, \hspace{0.1in} $S_2=\{T,R\}$
    \item Profiles: \hspace{0.45in} $S=S_1\times S_2=\{(T,T),(T,R),(R,T),(R,R)\}$
    \item Outcomes: \hspace{0.29in} $O=\{o_1,o_2,o_3,o_4\}$
    \item Outcome function: $f:S\to O$
\end{itemize}
\bsk
\centering
\begin{tabular}{c|c|c|c|c|}
\cline{2-5}
$s$: & $(T,T)$ & $(T,R)$ & $(R,T)$ & $(R,R)$ \\
\cline{2-5}
$f(s)$: & $o_1$ & $o_2$ & $o_3$ & $o_4$ \\
\cline{2-5}
\end{tabular}
\vfill
\raggedright
Notice: \textbf{nothing here says what anybody wants.}
\vspace{-0.4in}
\end{frame}

\begin{frame}{Definition: Game-Frame in Strategic Form}
A \textbf{game-frame in strategic form} is a list of four items
$$\langle I,\ (S_1,S_2,\ldots,S_n),\ O,\ f\rangle$$
where
\begin{itemize}
    \item $I=\{1,2,\ldots,n\}$ is the set of \textbf{players} ($n\ge 2$)
    \item $S_i$ is the set of \textbf{strategies} available to player $i$
    \begin{itemize}
        \item $S=S_1\times S_2\times\cdots\times S_n$ is the set of \textbf{strategy profiles}
        \item a typical profile is $s=(s_1,\ldots,s_n)$
    \end{itemize}
    \item $O$ is the set of \textbf{outcomes}
    \item $f:S\to O$ is the \textbf{outcome function}
\end{itemize}
\vfill
Why a function? \vspace{0.5in}
\end{frame}

\begin{frame}{Golden Balls as a Frame}
Write it out:
\bsk
\begin{itemize}
    \item $I=$ \vspace{0.3in}
    \item $S_1=$ \hspace{1.3in} $S_2=$ \vspace{0.3in}
    \item $O=$ \vspace{0.3in}
    \item $f=$ \vspace{0.3in}
\end{itemize}
\end{frame}

% ==============================================================
% 3. PREFERENCES  (extended treatment)
% ==============================================================

\begin{frame}[standout]
Preferences
\end{frame}

\begin{frame}{The Problem}
Outcomes are \textit{things}, not numbers:
\bsk
\begin{itemize}
    \item ``both at the Taco Truck''
    \item ``Sarah gets \$100,000 and Steven gets nothing''
    \item ``the park gets built and I pay \$40''
\end{itemize}
\bsk
We can't add them, average them, or take derivatives.
\vfill
All we will assume a player can do is \textbf{compare} them, two at a time.
\vspace{-0.5in}
\end{frame}

\begin{frame}{One Relation to Rule Them All}
For each player $i$ we write down a single relation on $O$:
\bsk
\centering
\large
$o \succsim_i o'$
\bsk
\normalsize
``Player $i$ considers $o$ to be \textbf{at least as good as} $o'$''
\bsk
\raggedright
\vfill
This one relation is enough. From it we \textit{define} the other two:
\bsk
\begin{itemize}
    \item \textbf{Strict preference:} $o\succ_i o'$ \ if \ $o\succsim_i o'$ \ and \ \textit{not} $o'\succsim_i o$
    \item \textbf{Indifference:} \hspace{0.16in} $o\sim_i o'$ \ if \ $o\succsim_i o'$ \ and \ $o'\succsim_i o$
\end{itemize}
\end{frame}

\begin{frame}{You've Done This Before}
You already know a relation that works exactly this way: $\ge$ on numbers.
\bsk
\centering
\begin{tabular}{|c|c|l|}
\hline
\textbf{Numbers} & \textbf{Outcomes} & \\
\hline
$x\ge y$ & $o\succsim_i o'$ & at least as good as \\
\hline
$x> y$ & $o\succ_i o'$ & strictly better than \\
\hline
$x= y$ & $o\sim_i o'$ & just as good as \\
\hline
\end{tabular}
\bsk
\raggedright
And the definitions are the same moves you'd make in grade school:
\begin{itemize}
    \item $x>y$ \ means \ $x\ge y$ \ and \ \textit{not} $y\ge x$
    \item $x=y$ \ means \ $x\ge y$ \ and \ $y\ge x$
\end{itemize}
\vspace{-0.3in}
\end{frame}

\begin{frame}{Where the Analogy Breaks}
With numbers, ``at least as big'' comes for free:
\begin{itemize}
    \item any two numbers \textit{can} be compared
    \item the comparisons \textit{never} cycle
\end{itemize}
\bsk
With outcomes, neither is automatic. ``Both at the taco truck'' and
``Sarah gets \$100,000'' are not on a number line.
\vfill
So the two properties we get for free with $\ge$ become
\textbf{assumptions} about $\succsim_i$.
\bsk
That's next.
\vspace{-0.4in}
\end{frame}

\begin{frame}{Reading the Symbols}
\centering
\begin{tabular}{|c|l|}
\hline
\textbf{Notation} & \textbf{Interpretation} \\
\hline
$o\succsim_i o'$ & $i$ finds $o$ at least as good as $o'$ \\
 & (better than, \textit{or} just as good as) \\
\hline
$o\succ_i o'$ & $i$ finds $o$ better than $o'$; $i$ \textit{prefers} $o$ to $o'$ \\
\hline
$o\sim_i o'$ & $i$ finds $o$ just as good as $o'$; $i$ is \textit{indifferent} \\
\hline
\end{tabular}
\bsk
\raggedright
The subscript $i$ matters: preferences belong to \textbf{a person}.
\bsk
$M\succ_{\text{Alice}} J$ \ but \ $M\sim_{\text{Bob}} J$ \ is perfectly consistent.
\vfill
Chains are read left to right, best to worst:
$$o_3\succ o_1\succ o_2\sim o_4$$
\vspace{-0.4in}
\end{frame}

\begin{frame}{Two Assumptions We Will Always Make}
\textbf{1. Completeness.} For any two outcomes $o$ and $o'$:
$$o\succsim_i o' \quad\text{or}\quad o'\succsim_i o \quad\text{(or both)}$$
\begin{itemize}
    \item i.e., the player is never unable to compare
    \item ``I don't know'' and ``I don't care'' are different! \vspace{0.4in}
\end{itemize}
\bsk
\textbf{2. Transitivity.} For any three outcomes:
$$\text{if } o_1\succsim_i o_2 \ \text{ and } \ o_2\succsim_i o_3 \ \text{ then } \ o_1\succsim_i o_3$$
\begin{itemize}
    \item i.e., the ranking doesn't cycle \vspace{0.4in}
\end{itemize}
\end{frame}

\begin{frame}{Why Transitivity?}
Suppose your preferences \textit{cycle}: \ $x\succ y$, \ $y\succ z$, \ but \ $z\succ x$.
\bsk
You own $z$. I own $x$ and $y$.
\bsk
\begin{itemize}
    \item You prefer $y$ to $z$, so you'd pay me a penny to trade $z$ for $y$ \vspace{0.35in}
    \item You prefer $x$ to $y$, so you'd pay me a penny to trade $y$ for $x$ \vspace{0.35in}
    \item You prefer $z$ to $x$, so you'd pay me a penny to trade $x$ for $z$ \vspace{0.35in}
\end{itemize}
Now you're back where you started, three cents poorer. Repeat forever.
\vfill
{\scriptsize This ``money pump'' argument is not in Bonanno; it's the standard defense of the transitivity assumption.}
\end{frame}

\begin{frame}{Transitivity Buys Us a Lot}
If $\succsim_i$ is complete and transitive, then $\succ_i$ and $\sim_i$ inherit good behavior:
\bsk
\begin{itemize}
    \item if $o_1\sim o_2$ and $o_2\sim o_3$ then $o_1\sim o_3$
    \item if $o_1\succ o_2$ and $o_2\succ o_3$ then $o_1\succ o_3$
    \item if $o_1\succ o_2$ and $o_2\sim o_3$ then $o_1\succ o_3$
    \item if $o_1\sim o_2$ and $o_2\succ o_3$ then $o_1\succ o_3$
\end{itemize}
\vfill
Together these let us line the outcomes up from best to worst,
with ties allowed.
\vspace{-0.5in}
\end{frame}

\begin{frame}{Three Ways to Write the Same Ranking}
Take $O=\{o_1,o_2,o_3,o_4,o_5\}$ and the ranking
$$o_3\ \succ\ o_5\sim o_1\ \succ\ o_4\ \succ\ o_2$$
\bsk
\textbf{Way 1: chain notation.} Exactly what's written above.
\bsk
\textbf{Way 2: a list, best on top.} Ties share a row.
\bsk
\centering
\begin{tabular}{rl}
best & $o_3$ \\
 & $o_1,\ o_5$ \\
 & $o_4$ \\
worst & $o_2$ \\
\end{tabular}
\vfill
\end{frame}

\begin{frame}{Three Ways to Write the Same Ranking}
\textbf{Way 3: attach a number to each outcome.}
\bsk
\centering
\begin{tabular}{c|c|c|c|c|c|}
\cline{2-6}
outcome & $o_1$ & $o_2$ & $o_3$ & $o_4$ & $o_5$ \\
\cline{2-6}
$U$ & $6$ & $1$ & $8$ & $2$ & $6$ \\
\cline{2-6}
\end{tabular}
\bsk
\raggedright
Rule: bigger number $=$ better; equal numbers $=$ indifferent.
\vfill
Check it against $o_3\succ o_5\sim o_1\succ o_4\succ o_2$: \vspace{0.7in}
\end{frame}

\begin{frame}{Definition: Ordinal Utility Function}
Let $\succsim$ be a complete and transitive ranking of a finite set $O$.
\bsk
A function $U:O\to\RR$ is an \textbf{ordinal utility function representing} $\succsim$ if,
for all outcomes $o$ and $o'$:
\bsk
\centering
$U(o)>U(o')$ \ if and only if \ $o\succ o'$
\bsk
$U(o)=U(o')$ \ if and only if \ $o\sim o'$
\bsk
\raggedright
\vfill
Equivalently: \hspace{0.3in} $o\succsim o'$ \ if and only if \ $U(o)\ge U(o')$.
\bsk
$U(o)$ is called the \textbf{utility} of outcome $o$.
\end{frame}

\begin{frame}{Utility Numbers Are Arbitrary}
All three of these represent the \textit{same} ranking $o_3\succ o_1\succ o_2\sim o_4$:
\bsk
\centering
\begin{tabular}{c|c|c|c|c|}
\cline{2-5}
 & $o_1$ & $o_2$ & $o_3$ & $o_4$ \\
\cline{2-5}
$f$ & $5$ & $2$ & $10$ & $2$ \\
\cline{2-5}
$g$ & $0.8$ & $0.7$ & $1$ & $0.7$ \\
\cline{2-5}
$h$ & $27$ & $1$ & $100$ & $1$ \\
\cline{2-5}
\end{tabular}
\bsk
\raggedright
There are infinitely many. Any order-preserving relabeling works.
\bsk
This is what ``\textbf{ordinal}'' means: \textit{only the order carries information}.
\vfill
\textbf{Consequence: utility cannot be compared across people.}
\bsk
Alice's utility for pizza is $5$. Bob's is $4$.
\begin{itemize}
    \item Does Alice like pizza more than Bob does? \hspace{0.2in} \textbf{No idea.}
    \item Bob could have written $400$ instead of $4$ and said the same thing
    \item Each person's numbers only speak about \textit{that person's own} ranking
\end{itemize}
\vspace{-0.3in}
\end{frame}

\begin{frame}{What Utility Does \textit{Not} Mean}
``Alice's utility for Mexican food is 10.''
\begin{itemize}
    \item Meaningless on its own.
\end{itemize}
\bsk
``Alice's utility for Mexican is 10 and for Japanese is 5.''
\begin{itemize}
    \item Meaningful --- it says she prefers Mexican.
    \item Does \textbf{not} say Mexican is ``twice as good.''
    \item We could have used $100$ and $-25$ and said exactly the same thing.
\end{itemize}
\vfill
Also: a utility of $1$ does not mean ``first choice.'' Low number, low rank.
\bsk
Think of it as an index of satisfaction --- but hold that thought loosely.
\vspace{-0.4in}
\end{frame}

\begin{frame}{Careful: Two Things That Look Alike}
\centering
\begin{tabular}{|l|l|}
\hline
\textbf{Ordinal utility (now)} & \textbf{Cardinal utility (Ch.\ 5)} \\
\hline
Only order matters & Differences matter too \\
\hline
Any increasing relabeling & Only positive affine \\
\ \ is an equally good $U$ & \ \ transformations $aU+b$ \\
\hline
Enough for dominance, & Needed for randomization \\
\ \ Nash equilibrium & \ \ and expected utility \\
\hline
\end{tabular}
\bsk
\raggedright
\vfill
For the next several weeks, we live entirely in the left column.
\vspace{-0.4in}
\end{frame}

% ==============================================================
% 4. FROM FRAME TO GAME
% ==============================================================

\begin{frame}{Definition: Ordinal Game in Strategic Form}
An \textbf{ordinal game in strategic form} is
$$\langle I,\ (S_1,\ldots,S_n),\ O,\ f,\ (\succsim_1,\ldots,\succsim_n)\rangle$$
where
\begin{itemize}
    \item $\langle I,(S_1,\ldots,S_n),O,f\rangle$ is a game-frame, and
    \item each $\succsim_i$ is a complete and transitive ranking of $O$
\end{itemize}
\vfill
The only new ingredient is the list of rankings, one per player.
\vspace{-0.5in}
\end{frame}

\begin{frame}{Payoff Functions}
Replace each ranking $\succsim_i$ with a utility function $U_i$ that represents it.
\bsk
Then define player $i$'s \textbf{payoff function} $\pi_i:S\to\RR$ by
$$\pi_i(s)\ =\ U_i\big(f(s)\big)$$
\bsk
\begin{itemize}
    \item $f(s)$: the outcome that profile $s$ produces
    \item $U_i(\cdot)$: how much $i$ likes that outcome
    \item $\pi_i(s)$: the two composed --- utility \textit{as a function of strategies}
\end{itemize}
\vfill
The triple $\langle I,(S_1,\ldots,S_n),(\pi_1,\ldots,\pi_n)\rangle$ is the
\textbf{reduced} ordinal game.\\
``Reduced'' because we've thrown away the names of the outcomes.
\vspace{-0.3in}
\end{frame}

\begin{frame}{The Pipeline}
\centering
\bsk
\begin{tikzpicture}[>=stealth, node distance=1.0cm,
   box/.style={draw, rounded corners=3pt, minimum height=0.8cm, minimum width=2.4cm, font=\small}]
  \node[box] (s) {strategy profile $s$};
  \node[box, right=1.5cm of s] (o) {outcome $f(s)$};
  \node[box, right=1.5cm of o] (u) {number $U_i(f(s))$};
  \draw[->, thick] (s) -- node[above, font=\scriptsize] {$f$} (o);
  \draw[->, thick] (o) -- node[above, font=\scriptsize] {$U_i$} (u);
  \draw[->, thick, dashed] (s.south) -- ++(0,-0.9) -| node[below, pos=0.25, font=\scriptsize] {$\pi_i$} (u.south);
\end{tikzpicture}
\bsk
\raggedright
\vfill
Everything from here on is written in terms of $\pi_i$.\\
But remember what's underneath it.
\vspace{-0.4in}
\end{frame}

\begin{frame}{Friday Night: Three Different Games, One Frame}
\textbf{Story A.} Both would rather be together; Alex likes tacos, Blair likes ramen.
\bsk
\begin{itemize}
    \item Alex: $o_1\succ_1 o_4\succ_1 o_2\sim_1 o_3$
    \item Blair: $o_4\succ_2 o_1\succ_2 o_2\sim_2 o_3$
\end{itemize}
\bsk
\centering
\begin{matrixgame}[top={Tacos, Ramen}, left={Tacos, Ramen}, player names={Alex, Blair}]
  \payoffmatrix{3,2 & 1,1 \\ 1,1 & 2,3}
\end{matrixgame}
\vfill
\end{frame}

\begin{frame}{Friday Night: Three Different Games, One Frame}
\textbf{Story B.} Alex wants company; Blair wants to eat alone tonight.
\bsk
\begin{itemize}
    \item Alex: $o_1\sim_1 o_4\succ_1 o_2\sim_1 o_3$
    \item Blair: $o_2\sim_2 o_3\succ_2 o_1\sim_2 o_4$
\end{itemize}
\bsk
\centering
\begin{matrixgame}[top={Tacos, Ramen}, left={Tacos, Ramen}, player names={Alex, Blair}]
  \payoffmatrix{2,0 & 0,2 \\ 0,2 & 2,0}
\end{matrixgame}
\vfill
\raggedright
Same frame. Completely different game. \vspace{0.4in}
\end{frame}

\begin{frame}{Friday Night: Three Different Games, One Frame}
\textbf{Story C.} Neither cares about company at all; each just wants their favorite food.
\bsk
\begin{itemize}
    \item Alex: $o_1\sim_1 o_2\succ_1 o_3\sim_1 o_4$ \hspace{0.3in} (tacos, period)
    \item Blair: $o_2\sim_2 o_4\succ_2 o_1\sim_2 o_3$ \hspace{0.3in} (ramen, period)
\end{itemize}
\bsk
\centering
\begin{matrixgame}[top={Tacos, Ramen}, left={Tacos, Ramen}, player names={Alex, Blair}]
  \payoffmatrix{1,0 & 1,1 \\ 0,0 & 0,1}
\end{matrixgame}
\vfill
\raggedright
Here the answer feels obvious. Why? \vspace{0.5in}
\end{frame}

\begin{frame}{Back to Golden Balls}
Suppose \textbf{both} are selfish and greedy:
\begin{itemize}
    \item Sarah: $o_3\succ_1 o_1\succ_1 o_2\sim_1 o_4$
    \item Steven: $o_2\succ_2 o_1\succ_2 o_3\sim_2 o_4$
\end{itemize}
\bsk
Pick utilities (any numbers with the right order will do):
\bsk
\centering
\begin{tabular}{c|c|c|c|c|}
\cline{2-5}
 & $o_1$ & $o_2$ & $o_3$ & $o_4$ \\
\cline{2-5}
$U_1$ & $3$ & $2$ & $4$ & $2$ \\
\cline{2-5}
$U_2$ & $3$ & $4$ & $2$ & $2$ \\
\cline{2-5}
\end{tabular}
\bsk
\begin{matrixgame}[top={Split, Steal}, left={Split, Steal}, player names={Sarah, Steven}]
  \payoffmatrix{3,3 & 2,4 \\ 4,2 & 2,2}
\end{matrixgame}
\end{frame}

\begin{frame}{Back to Golden Balls}
Now suppose Sarah is \textbf{fair-minded and benevolent}, Steven still selfish and greedy:
\begin{itemize}
    \item Sarah: $o_1\succ_1 o_3\succ_1 o_2\succ_1 o_4$
    \item Steven: $o_2\succ_2 o_1\succ_2 o_3\sim_2 o_4$
\end{itemize}
\bsk
\centering
\begin{tabular}{c|c|c|c|c|}
\cline{2-5}
 & $o_1$ & $o_2$ & $o_3$ & $o_4$ \\
\cline{2-5}
$U_1$ & $4$ & $2$ & $3$ & $1$ \\
\cline{2-5}
$U_2$ & $3$ & $4$ & $2$ & $2$ \\
\cline{2-5}
\end{tabular}
\bsk
\begin{matrixgame}[top={Split, Steal}, left={Split, Steal}, player names={Sarah, Steven}]
  \payoffmatrix{4,3 & 2,4 \\ 3,2 & 1,2}
\end{matrixgame}
\vfill
\raggedright
Same show, same rules, same table of prizes. \textbf{Different game.}
\end{frame}

% ==============================================================
% 5. DOMINANCE
% ==============================================================

\begin{frame}[standout]
Dominance
\end{frame}

\begin{frame}{The Idea}
Fix one player. Compare two of \textit{her own} strategies, $a$ and $b$.
\bsk
Ask: how do $a$ and $b$ compare \textbf{in every situation she might face}?
\bsk
\begin{itemize}
    \item If $a$ always beats $b$: \hspace{0.42in} $a$ \textbf{strictly dominates} $b$
    \item If $a$ never loses to $b$, \\ and sometimes beats it: \hspace{0.06in} $a$ \textbf{weakly dominates} $b$
    \item If they always tie: \hspace{0.5in} $a$ and $b$ are \textbf{equivalent}
\end{itemize}
\vfill
``Situation'' $=$ a choice of strategies by \textit{everyone else}.
\vspace{-0.4in}
\end{frame}

\begin{frame}{Notation for ``Everyone Else''}
With $n$ players, a profile is $s=(s_1,\ldots,s_n)$.
\bsk
Fix a player $i$. Write
$$s_{-i}=\text{the strategies of all players other than } i$$
so that $s=(s_i,\ s_{-i})$.
\bsk
\begin{itemize}
    \item $S_{-i}$ is the set of all such sub-profiles: $S_{-i}=\times_{j\ne i}S_j$
    \item $\pi_i(a,s_{-i})$ means: $i$ plays $a$, everyone else plays $s_{-i}$
\end{itemize}
\vfill
Example: $S_1=\{a,b,c\}$, $S_2=\{d,e\}$, $S_3=\{f,g\}$, and $s=(b,d,g)$.
\bsk
Then $s_{-2}=$ \underline{\hspace{1.2in}} \hspace{0.3in} and $(e,s_{-2})=$ \underline{\hspace{1.2in}}
\end{frame}

\begin{frame}{An Example (Player 1's Payoffs Only)}
\centering
\begin{matrixgame}[top={L, C, R}, left={T, M, B}, player names={Player 1, Player 2}]
  \payoffmatrix{4,\cdot & 3,\cdot & 5,\cdot \\ 2,\cdot & 1,\cdot & 4,\cdot \\ 4,\cdot & 1,\cdot & 5,\cdot}
\end{matrixgame}
\bsk
\raggedright
Compare $T$ and $M$, row by row:
\begin{itemize}
    \item vs.\ $L$: \hspace{0.3in} $4$ \underline{\hspace{0.3in}} $2$
    \item vs.\ $C$: \hspace{0.3in} $3$ \underline{\hspace{0.3in}} $1$
    \item vs.\ $R$: \hspace{0.3in} $5$ \underline{\hspace{0.3in}} $4$
\end{itemize}
\vfill
Conclusion: \vspace{0.4in}
\end{frame}

\begin{frame}{An Example (Player 1's Payoffs Only)}
\centering
\begin{matrixgame}[top={L, C, R}, left={T, M, B}, player names={Player 1, Player 2}]
  \payoffmatrix{4,\cdot & 3,\cdot & 5,\cdot \\ 2,\cdot & 1,\cdot & 4,\cdot \\ 4,\cdot & 1,\cdot & 5,\cdot}
\end{matrixgame}
\bsk
\raggedright
Now compare $T$ and $B$:
\begin{itemize}
    \item vs.\ $L$: \hspace{0.3in} $4$ \underline{\hspace{0.3in}} $4$
    \item vs.\ $C$: \hspace{0.3in} $3$ \underline{\hspace{0.3in}} $1$
    \item vs.\ $R$: \hspace{0.3in} $5$ \underline{\hspace{0.3in}} $5$
\end{itemize}
\vfill
$T$ never does worse, and does better against $C$. This is \textbf{weak} dominance.
\vspace{-0.3in}
\end{frame}

\begin{frame}{Definition: Dominance}
Let $i$ be a player and $a,b\in S_i$ two of her strategies.
\bsk
\begin{itemize}
    \item $a$ \textbf{strictly dominates} $b$ if
    $$\pi_i(a,s_{-i})>\pi_i(b,s_{-i}) \quad\text{for every } s_{-i}\in S_{-i}$$
    \bsk
    \item $a$ \textbf{weakly dominates} $b$ if
    $$\pi_i(a,s_{-i})\ge\pi_i(b,s_{-i}) \quad\text{for every } s_{-i}\in S_{-i}$$
    \hspace{0.2in} \textit{and} \ $\pi_i(a,s_{-i})>\pi_i(b,s_{-i})$ for at least one $s_{-i}$
    \bsk
    \item $a$ is \textbf{equivalent} to $b$ if
    $$\pi_i(a,s_{-i})=\pi_i(b,s_{-i}) \quad\text{for every } s_{-i}\in S_{-i}$$
\end{itemize}
\end{frame}

\begin{frame}{Two Conventions}
\textbf{1.} Strict dominance technically satisfies the definition of weak dominance.
\bsk
From here on, ``$a$ weakly dominates $b$'' means
\textit{weakly but not strictly}.
\vfill
\textbf{2.} ``Dominated'' is a comparison, so name the other strategy.
\bsk
\begin{itemize}
    \item ``$x$ is dominated'' \hspace{0.35in} $\rightsquigarrow$ \hspace{0.1in} dominated by \textit{what}?
    \item ``$x$ is dominated by $y$'' \hspace{0.05in} $\rightsquigarrow$ \hspace{0.1in} fine
\end{itemize}
\bsk
Think of ``dominates'' as ``is better than.''
\vspace{-0.4in}
\end{frame}

\begin{frame}{Dominant $=$ Best}
``Dominant'' needs no second strategy named --- it means \textit{better than all of them}.
\bsk
\begin{itemize}
    \item $a$ is a \textbf{strictly dominant strategy} if $a$ strictly dominates
      \textit{every} other strategy of player $i$
    \bsk
    \item $a$ is a \textbf{weakly dominant strategy} if, for every other strategy $x$,
      either $a$ weakly dominates $x$ or $a$ is equivalent to $x$
\end{itemize}
\vfill
Consequences:
\begin{itemize}
    \item There is at most \textbf{one} strictly dominant strategy. Why? \vspace{0.35in}
    \item There can be \textbf{several} weakly dominant strategies --- but any two of
      them must be equivalent.
\end{itemize}
\end{frame}

\begin{frame}{Dominant-Strategy Profiles}
Let $s=(s_1,\ldots,s_n)$ be a strategy profile.
\bsk
\begin{itemize}
    \item $s$ is a \textbf{strict dominant-strategy profile} if $s_i$ is strictly dominant
      for \textit{every} player $i$
    \bsk
    \item $s$ is a \textbf{weak dominant-strategy profile} if $s_i$ is weakly dominant for
      every player $i$, and for at least one player $j$, $s_j$ is not strictly dominant
\end{itemize}
\vfill
Unqualified ``dominant-strategy profile'' defaults to \textbf{weak}.
\bsk
{\footnotesize (You'll also see ``dominant-strategy equilibrium'' in the literature.)}
\vspace{-0.4in}
\end{frame}

\begin{frame}{Your Turn}
Only Player 1's payoffs are shown.
\bsk
\centering
\begin{matrixgame}[top={E, F, G}, left={A, B, C, D}, player names={Player 1, Player 2}]
  \payoffmatrix{3,\cdot & 2,\cdot & 1,\cdot \\ 2,\cdot & 1,\cdot & 0,\cdot \\ 3,\cdot & 2,\cdot & 1,\cdot \\ 2,\cdot & 0,\cdot & 0,\cdot}
\end{matrixgame}

\raggedright
Classify each pair. Then: any weakly dominant strategies?
\bsk
\begin{tabular}{ll}
$A$ vs.\ $B$: \hspace{1.4in} & $A$ vs.\ $D$: \\[0.3in]
$A$ vs.\ $C$: & $B$ vs.\ $D$: \\[0.3in]
\end{tabular}
\end{frame}

% ==============================================================
% 6. GOLDEN BALLS PAYOFF + THE VIDEO
% ==============================================================

\begin{frame}{So What About Sarah?}
Selfish and greedy version:
\bsk
\centering
\begin{matrixgame}[top={Split, Steal}, left={Split, Steal}, player names={Sarah, Steven}]
  \payoffmatrix{3,3 & 2,4 \\ 4,2 & 2,2}
\end{matrixgame}
\bsk
\raggedright
\begin{itemize}
    \item Steal is a \underline{\hspace{1.2in}} dominant strategy for each player
    \item So (Steal, Steal) is a \underline{\hspace{1.2in}} dominant-strategy profile
\end{itemize}
\vfill
And the fair-minded version? \vspace{0.5in}
\end{frame}

\begin{frame}{Let's Watch}
\vfill
\centering
\large
\textbf{Golden Balls: Split or Steal}
\bsk
\normalsize
\url{https://www.youtube.com/watch?v=tBtr8-VMj0E}
\vfill
\raggedright
While you watch, ask yourself:
\begin{itemize}
    \item Which game are \textit{they} playing?
    \item What does the talking do, if agreements aren't binding?
\end{itemize}
\vspace{-0.3in}
\end{frame}

% ==============================================================
% 7. PRISONER'S DILEMMA & PARETO
% ==============================================================

\begin{frame}{Our Own Example: The Price War}
Two food trucks park on the same block every day. Each morning, each owner
independently sets a \textbf{Normal} price or runs a \textbf{Discount}.
\bsk
\begin{itemize}
    \item If neither discounts: they split the lunch crowd at healthy margins
    \item If exactly one discounts: that truck pulls most of the crowd
    \item If both discount: they split the same crowd at thin margins
\end{itemize}
\bsk
Each owner ranks the four outcomes:
$$\text{(I discount, they don't)}\ \succ\ \text{(neither)}\ \succ\ \text{(both)}\ \succ\ \text{(they discount, I don't)}$$
\vspace{-0.4in}
\end{frame}

\begin{frame}{The Price War}
\centering
\begin{matrixgame}[top={Normal, Discount}, left={Normal, Discount}, player names={Truck 1, Truck 2}, br]
  \payoffmatrix{2,2 & 0,3 \\ 3,0 & 1,1}
\end{matrixgame}
\bsk
\raggedright
\begin{itemize}
    \item Discount \underline{\hspace{1.4in}} dominates Normal, for both
    \item So (Discount, Discount) is a \underline{\hspace{1.4in}} profile
\end{itemize}
\vfill
But look at (Normal, Normal): \vspace{0.6in}
\end{frame}

\begin{frame}{Definition: Pareto Superiority}
Let $o$ and $o'$ be two outcomes.
\bsk
\begin{itemize}
    \item $o$ is \textbf{strictly Pareto superior} to $o'$ if \textit{every} player prefers
      $o$ to $o'$: \ $o\succ_i o'$ for all $i$
    \bsk
    \item $o$ is \textbf{weakly Pareto superior} to $o'$ if every player finds $o$ at least
      as good, and at least one strictly prefers it
\end{itemize}
\bsk
In reduced games, same thing with payoffs: $\pi_i(s)>\pi_i(s')$ for all $i$, etc.
\vfill
In the Price War, (Normal, Normal) is strictly Pareto superior to (Discount, Discount).
\bsk
Yet rational play lands on (Discount, Discount).
\vspace{-0.3in}
\end{frame}

\begin{frame}{The Prisoner's Dilemma}
That structure has a name. Bonanno's version:
\bsk
Doug and Ed are the only two candidates for a best-worker prize, currently tied.
Unless one outperforms the other, nobody gets it. Each chooses Normal or Extra effort.
\bsk
\begin{itemize}
    \item $o_1$: nobody gets the prize, nobody sacrifices family time
    \item $o_2$: Ed gets the prize and sacrifices family time, Doug does not
    \item $o_3$: Doug gets the prize and sacrifices family time, Ed does not
    \item $o_4$: nobody gets the prize, both sacrifice family time
\end{itemize}
\bsk
Both are willing to sacrifice family time for the prize, otherwise value it, and are
envious: \ $o_3\succ_{\text{Doug}} o_1\succ_{\text{Doug}} o_4\succ_{\text{Doug}} o_2$.
\vspace{-0.3in}
\end{frame}

\begin{frame}{The Prisoner's Dilemma}
\centering
\begin{matrixgame}[top={Normal, Extra}, left={Normal, Extra}, player names={Doug, Ed}, br]
  \payoffmatrix{2,2 & 0,3 \\ 3,0 & 1,1}
\end{matrixgame}
\bsk
\raggedright
Extra effort is strictly dominant for both. \\
(Extra, Extra) is a strict dominant-strategy profile.
\vfill
Same game as the price war --- different story, identical structure.\\
That is what game theory is for.
\vspace{-0.4in}
\end{frame}

\begin{frame}{Individual vs.\ Collective Rationality}
A strictly dominant strategy is not a suggestion. Playing anything else means a
lower payoff \textit{no matter what the other player does}.
\bsk
So rational individual play $\Rightarrow$ (Extra, Extra), even though both prefer (Normal, Normal).
\vfill
Couldn't they just agree?
\begin{itemize}
    \item Non-cooperative game theory assumes agreements are \textbf{not binding}
    \item If I expect you to keep it, I gain by breaking it \vspace{0.3in}
    \item If I expect you to break it, I gain by breaking it too \vspace{0.3in}
\end{itemize}
\end{frame}

\begin{frame}{Where We Landed}
\begin{itemize}
    \item A \textbf{game-frame} is players, strategies, outcomes, and $f$
    \item Preferences are a complete, transitive relation $\succsim_i$ on outcomes
    \item An \textbf{ordinal utility function} represents $\succsim_i$; the numbers are arbitrary
    \item Frame $+$ preferences $=$ \textbf{game}; composing gives payoffs $\pi_i=U_i\circ f$
    \item \textbf{Dominance} compares two strategies of one player across all situations
    \item A dominant strategy is best against everything
\end{itemize}
\vfill
Next: what if nobody has a dominant strategy?
\vspace{-0.5in}
\end{frame}

\begin{frame}{Reading \& Practice}
\begin{itemize}
    \item Bonanno, \textit{Game Theory}, \S2.1--2.2
    \item Exercises: \S2.9.1 (frames and preferences), \S2.9.2 (dominance)
    \item Full solutions are in \S2.10 --- try them before you look
\end{itemize}
\vfill
\end{frame}

\end{document}
